This section defines the mechanistic framework used to analyze how optimizer
choice can affect feature-based uncertainty. We do not introduce a new OOD
detector. Instead, we treat optimizer choice as an intervention on penultimate
representation geometry, and we treat each feature-based detector family as a
readout of a particular geometry channel. The resulting question is not which
optimizer is universally better, but which optimizer-induced geometry is
compatible with which downstream readout.

% TODO: Optional method schematic: optimizer -> geometry fingerprint -> detector readout -> score ranking.

Table~\ref{tab:readout-map} summarizes the detector families considered by this
framework. The table is conceptual: it records which geometry channel each
readout can be sensitive to and which diagnostic control helps localize that
channel.

\begin{table}[H]
\caption{Feature-based detector families as geometry readouts. The entries
summarize the framework only; they do not report experimental results.}
\label{tab:readout-map}
\begin{center}
\footnotesize
\setlength{\tabcolsep}{3pt}
\renewcommand{\arraystretch}{1.15}
\begin{tabular}{@{}>{\raggedright\arraybackslash}p{0.17\textwidth}>{\raggedright\arraybackslash}p{0.25\textwidth}>{\raggedright\arraybackslash}p{0.27\textwidth}>{\raggedright\arraybackslash}p{0.20\textwidth}@{}}
\toprule
Detector family & Score form & Geometry channel & Diagnostic control \\
\midrule
Mahalanobis
& \(-\min_c (f-\mu_c)^\top \Sigma_W^{-1}(f-\mu_c)\)
& class centers, within-class covariance, precision spectrum
& post-hoc L2 normalization; covariance shrinkage \\
Gaussian density
& \(\log\sum_c \pi_c \mathcal{N}(f\mid\mu_c,\widehat{\Sigma}_c)\)
& Gaussian density, covariance volume, anisotropy, conditioning
& tied, diagonal, and shrinkage covariance \\
kNN
& \(-d_k(f,\mathcal{B})\)
& local feature-bank spacing, radial scale, angular neighborhoods
& raw kNN vs. L2-kNN \\
CTM angular readout
& \(\max_c \langle \widehat f,a_c\rangle\)
& prototype or classifier directions, angular class separation
& normalized features; prototype vs. classifier direction \\
NECO collapse/residual readout
& \(\max_c\langle\widehat f,\widehat\mu_c\rangle\), \(-\|(I-P_M)f\|_2^2\)
& collapse-like prototype geometry, residual energy outside class-mean subspace
& prototype vs. residual readout; subspace dimension \\
\bottomrule
\end{tabular}
\end{center}
\end{table}


\subsection{Optimizer-conditioned representation geometry}
\label{subsec:optimizer-conditioned-geometry}

Let \(o\) denote an optimizer or training rule. A classifier trained under
optimizer \(o\) is written as
\begin{equation}
    z_o(x)=W_o f_o(x)+b_o,
    \label{eq:classifier-decomposition}
\end{equation}
where \(f_o(x)\in\mathbb{R}^d\) is the penultimate feature,
\(W_o\in\mathbb{R}^{K\times d}\) is the classifier weight matrix, and
\(z_o(x)\in\mathbb{R}^K\) denotes the logits. Let
\(\mathcal{D}_{\mathrm{tr}}\) be the ID training set and
\(\mathcal{D}_c\subset\mathcal{D}_{\mathrm{tr}}\) be the subset of class \(c\).
For each optimizer \(o\), we define the class mean and covariance in the
penultimate feature space as
\begin{equation}
    \mu_{c,o}
    =
    \frac{1}{|\mathcal{D}_c|}
    \sum_{(x_i,y_i)\in\mathcal{D}_c}
    f_o(x_i),
    \label{eq:class-mean}
\end{equation}
\begin{equation}
    \Sigma_{c,o}
    =
    \frac{1}{|\mathcal{D}_c|-1}
    \sum_{(x_i,y_i)\in\mathcal{D}_c}
    \bigl(f_o(x_i)-\mu_{c,o}\bigr)
    \bigl(f_o(x_i)-\mu_{c,o}\bigr)^\top.
    \label{eq:class-covariance}
\end{equation}
When a detector uses a shared covariance estimate, we use the pooled
within-class covariance
\begin{equation}
    \Sigma_{W,o}
    =
    \frac{1}{K}
    \sum_{c=1}^{K}
    \Sigma_{c,o}.
    \label{eq:pooled-within-covariance}
\end{equation}
We also define the ID feature bank
\begin{equation}
    \mathcal{B}_o
    =
    \{f_o(x_i):(x_i,y_i)\in\mathcal{D}_{\mathrm{tr}}\},
    \label{eq:id-feature-bank}
\end{equation}
which is used by neighborhood-based detectors.

To obtain a controlled optimizer-side geometry intervention, we use an
AdamW-to-Adam interpolation along the weight-decay coupling axis. For a fixed
total weight decay \(\lambda_{\mathrm{total}}\), define
\begin{equation}
    \lambda_{\mathrm{coupled}}
    =
    \gamma\lambda_{\mathrm{total}},
    \qquad
    \lambda_{\mathrm{decoupled}}
    =
    (1-\gamma)\lambda_{\mathrm{total}},
    \qquad
    \gamma\in[0,1].
    \label{eq:adamw-adam-interpolation}
\end{equation}
The coupled component is included in the gradient update, whereas the
decoupled component is applied as a separate parameter shrinkage. Abstractly,
for an optimizer family with preconditioning or momentum operator
\(\mathcal{P}_o\), the update can be written as
\begin{equation}
    g_t^{(\gamma)}
    =
    \nabla_\theta \mathcal{L}(\theta_t)
    +
    \lambda_{\mathrm{coupled}}\theta_t,
    \label{eq:optimizer-coupled-gradient}
\end{equation}
\begin{equation}
    \theta_{t+1}
    =
    (1-\eta_t\lambda_{\mathrm{decoupled}})\theta_t
    -
    \eta_t\mathcal{P}_o(g_t^{(\gamma)}).
    \label{eq:optimizer-interpolation-update}
\end{equation}
In the Adam family, \(\gamma=0\) corresponds to AdamW decoupled weight
decay and \(\gamma=1\) corresponds to Adam coupled L2 regularization. We
use this interpolation not to claim that all optimizer effects reduce to a
single scalar, but to create a controlled path along which representation
geometry can be compared while the detector pipeline is fixed.

For each trained model, we summarize the penultimate representation by a
geometry fingerprint
\begin{equation}
    G_o
    =
    \bigl[
    G_o^{\mathrm{norm}},
    G_o^{\mathrm{cov}},
    G_o^{\mathrm{NC}},
    G_o^{\mathrm{ang/sub}}
    \bigr].
    \label{eq:geometry-fingerprint}
\end{equation}
The norm component includes feature-norm statistics; the covariance component
includes within-class dispersion, covariance spectrum, anisotropy,
conditioning, and effective rank; the neural-collapse component includes
neural-collapse measures of within-class collapse, class-mean geometry, and
classifier--feature alignment; and the angular/subspace component describes
prototype directions and residual feature directions. The role of
Section~\ref{sec:experiments} is to test how these geometry components move
under optimizer interventions and how detector families respond.

\subsection{AUROC ranking and optimizer-induced score shifts}
\label{subsec:ranking-score-shifts}

All detector scores are oriented so that larger values indicate more ID-like
inputs. For a score \(S\), OOD AUROC can be written as the probability that a
random ID example receives a higher score than a random OOD example:
\begin{equation}
    \mathrm{AUROC}(S)
    =
    \mathbb{P}\!\left[
    S(X_{\mathrm{ID}})
    >
    S(X_{\mathrm{OOD}})
    \right]
    +
    \frac{1}{2}
    \mathbb{P}\!\left[
    S(X_{\mathrm{ID}})
    =
    S(X_{\mathrm{OOD}})
    \right].
    \label{eq:auroc-ranking}
\end{equation}
Thus, AUROC depends on ID/OOD score orderings rather than the absolute scale of
the score. An optimizer-induced change in score magnitude is therefore not by
itself a detector failure. It becomes detector-relevant only when it changes
enough pairwise ID/OOD rankings.

Let \(m\) index a detector family. Relative to a reference optimizer
\(o_{\mathrm{ref}}\), we write the score under optimizer \(o\) as
\begin{equation}
    S_m^{(o)}(x)
    =
    S_m^{(o_{\mathrm{ref}})}(x)
    +
    \Delta_m^{(o)}(x),
    \label{eq:optimizer-induced-score-shift}
\end{equation}
where \(\Delta_m^{(o)}(x)\) is the sample-dependent score shift induced by the
optimizer-conditioned geometry. If the score shift were equivalent to a
strictly increasing transformation of the reference score, the ranking would be
preserved. In contrast, geometry changes can act differently on different
samples. For an ID sample \(x\) and an OOD sample \(x'\) that are correctly
ordered under the reference optimizer, a ranking flip occurs when
\begin{equation}
    \Delta_m^{(o)}(x')
    -
    \Delta_m^{(o)}(x)
    >
    S_m^{(o_{\mathrm{ref}})}(x)
    -
    S_m^{(o_{\mathrm{ref}})}(x').
    \label{eq:ranking-flip-condition}
\end{equation}
This inequality is not a theorem or an additional detector objective. It only
formalizes the ranking perspective used throughout the paper. The substantive
question is how each feature-based detector produces the sample-dependent shift
\(\Delta_m^{(o)}(x)\) from the geometry it reads.

\subsection{Precision and density readouts}
\label{subsec:precision-density-readouts}

\paragraph{Mahalanobis readout.}
Mahalanobis scoring reads class centers through a precision-weighted distance.
Under optimizer \(o\), the score is
\begin{equation}
    S_{\mathrm{Maha}}^{(o)}(x)
    =
    -
    \min_{c}
    \bigl(f_o(x)-\mu_{c,o}\bigr)^\top
    \Sigma_{W,o}^{-1}
    \bigl(f_o(x)-\mu_{c,o}\bigr).
    \label{eq:mahalanobis-score}
\end{equation}
This score depends on the location of class means, the spread of within-class
features, and the covariance spectrum that determines the precision matrix
\(\Sigma_{W,o}^{-1}\). If an optimizer changes within-class dispersion or
expands particular covariance directions, the precision-weighted penalty can
change along those directions. In the simplified case where the feature vector
and class means are fixed and
\(\Sigma_{W,o}\succeq\Sigma_{W,o_{\mathrm{ref}}}\), matrix inversion reverses
the positive-semidefinite order:
\begin{equation}
    \Sigma_{W,o}^{-1}
    \preceq
    \Sigma_{W,o_{\mathrm{ref}}}^{-1}.
    \label{eq:precision-order-reversal}
\end{equation}
Consequently, for any fixed residual \(f-\mu_c\),
\begin{equation}
    (f-\mu_c)^\top
    \Sigma_{W,o}^{-1}
    (f-\mu_c)
    \le
    (f-\mu_c)^\top
    \Sigma_{W,o_{\mathrm{ref}}}^{-1}
    (f-\mu_c).
    \label{eq:precision-quadratic-order}
\end{equation}
This condition does not imply that Mahalanobis AUROC must improve or degrade.
Rather, it identifies one detector-relevant channel: covariance changes can
make samples appear closer or farther from class means under the
precision-weighted metric. Because this shift is direction-dependent, samples
aligned with affected covariance directions can receive different
\(\Delta_m^{(o)}(x)\) values and can therefore alter rankings.

\paragraph{Gaussian density readout.}
We use Gaussian feature-density scores as post-hoc readouts of class
density. The generic score is
\begin{equation}
    S_{\mathrm{GMM}}^{(o)}(x)
    =
    \log
    \sum_{c=1}^{K}
    \pi_c
    \mathcal{N}
    \bigl(
    f_o(x)
    \mid
    \mu_{c,o},
    \widehat{\Sigma}_{c,o}
    \bigr).
    \label{eq:gmm-score}
\end{equation}
For a fixed class component, the log density decomposes into a quadratic term
and a covariance-volume term:
\begin{equation}
    \log
    \mathcal{N}
    (f\mid \mu_c,\Sigma_c)
    =
    -
    \frac{1}{2}
    (f-\mu_c)^\top
    \Sigma_c^{-1}
    (f-\mu_c)
    -
    \frac{1}{2}
    \log|\Sigma_c|
    +
    \mathrm{const}.
    \label{eq:gaussian-log-density-decomposition}
\end{equation}
The first term is a precision-weighted quadratic penalty, as in Mahalanobis
scoring. The second term penalizes covariance volume. Gaussian density scores
are therefore sensitive not only to how far a feature is from a class mean, but
also to the estimated volume and conditioning of each class covariance.

This makes the density readout richer, but also more delicate. Covariance
expansion can reduce the quadratic penalty for some directions, while the
log-determinant term penalizes larger covariance volume. Since the mixture
score combines class components through a log-sum-exp, changes in covariance
scale, anisotropy, class overlap, and component selection can all affect
ID/OOD rankings. We therefore treat tied, diagonal, and shrinkage covariance
variants as diagnostic density readouts rather than as a claim that a single
Gaussian estimator fully captures the representation.

\subsection{Neighborhood readouts}
\label{subsec:neighborhood-readouts}

kNN scoring reads the local spacing of the ID feature bank. Under optimizer
\(o\), we define
\begin{equation}
    S_{\mathrm{kNN}}^{(o)}(x)
    =
    -
    d_k
    \bigl(
    f_o(x),
    \mathcal{B}_o
    \bigr),
    \label{eq:knn-score}
\end{equation}
where \(d_k(f,\mathcal{B}_o)\) is the Euclidean distance from \(f\) to its
\(k\)-th nearest neighbor in the ID training feature bank. Unlike Mahalanobis
and Gaussian density scores, kNN does not explicitly fit a covariance matrix.
However, it remains sensitive to representation geometry because the ID feature
bank itself changes with the optimizer.

The raw Euclidean distance between a query feature \(f=ru\) and a bank feature
\(f_i=r_i u_i\), with \(\|u\|_2=\|u_i\|_2=1\), satisfies
\begin{equation}
    \|f-f_i\|_2^2
    =
    r^2+r_i^2-2rr_i\langle u,u_i\rangle.
    \label{eq:euclidean-radial-angular-decomposition}
\end{equation}
Thus, raw kNN mixes radial scale, angular similarity, and local feature
density. If an optimizer changes feature-norm distributions or local cluster
spacing, raw kNN scores can shift even when angular neighborhoods are similar.
After L2 normalization,
\begin{equation}
    \|\widehat f-\widehat f_i\|_2^2
    =
    2
    -
    2\langle \widehat f,\widehat f_i\rangle,
    \qquad
    \widehat f=\frac{f}{\|f\|_2+\epsilon},
    \label{eq:l2-angular-distance}
\end{equation}
so the distance becomes angular up to the normalization constant. Comparing raw
kNN and L2-kNN therefore tests whether neighborhood-based detector behavior is
driven by local density itself or by radial scale coupled to the feature bank.

\subsection{Angular and collapse-based readouts}
\label{subsec:angular-collapse-readouts}

Precision, density, and neighborhood scores read distances or densities in raw
feature space. Other feature-based detectors read angular prototype structure
or neural-collapse structure. We include these readouts because optimizer
choice can change class-mean directions, classifier alignment, and residual
feature subspaces even when ID accuracy remains similar.

\paragraph{CTM angular readout.}
CTM scores read cosine similarity between a query feature and
class-related directions. A generic angular readout can be written as
\begin{equation}
    S_{\mathrm{cos}}^{(o)}(x)
    =
    \max_c
    \left\langle
    \frac{f_o(x)}{\|f_o(x)\|_2+\epsilon},
    a_{c,o}
    \right\rangle,
    \qquad
    \|a_{c,o}\|_2=1,
    \label{eq:ctm-cosine-readout}
\end{equation}
where \(a_{c,o}\) may be a normalized class prototype, a normalized classifier
weight direction, or a detector-specific class direction. Because the score
uses normalized features, it is less sensitive to global feature norm than raw
Mahalanobis or raw kNN. It is instead sensitive to angular class separation,
prototype stability, and classifier--feature alignment. If an optimizer changes
the directions of class means or classifier weights, or weakens their
alignment, angular readouts can change even when radial diagnostics show little
movement.

\paragraph{NECO collapse and residual readout.}
NECO scores use neural-collapse structure to define prototype and
residual signals. Let the centered class-mean matrix be
\begin{equation}
    M_o
    =
    \bigl[
    \widetilde{\mu}_{1,o},
    \ldots,
    \widetilde{\mu}_{K,o}
    \bigr],
    \qquad
    \widetilde{\mu}_{c,o}
    =
    \mu_{c,o}
    -
    \frac{1}{K}
    \sum_{j=1}^{K}
    \mu_{j,o}.
    \label{eq:centered-class-mean-matrix}
\end{equation}
A prototype-alignment readout can be expressed as
\begin{equation}
    S_{\mathrm{proto}}^{(o)}(x)
    =
    \max_c
    \left\langle
    \frac{f_o(x)}{\|f_o(x)\|_2+\epsilon},
    \frac{\widetilde{\mu}_{c,o}}
    {\|\widetilde{\mu}_{c,o}\|_2+\epsilon}
    \right\rangle.
    \label{eq:neco-prototype-readout}
\end{equation}
A residual-subspace readout can be written as
\begin{equation}
    S_{\mathrm{res}}^{(o)}(x)
    =
    -
    \left\|
    (I-P_{M_o}) f_o(x)
    \right\|_2^2,
    \label{eq:neco-residual-readout}
\end{equation}
where \(P_{M_o}\) is the projection onto the span of the centered class means.
The exact NECO implementation can combine these ingredients differently,
but the relevant geometry channels are the same: within-class collapse,
simplex-like class-mean structure, classifier--prototype self-duality, and
residual energy outside the class-mean subspace.

This readout family is important for our framework because it is not primarily
a covariance estimator. It asks whether the representation has the
collapse-like prototype geometry that prototype and residual scores require.
Optimizers that preserve accuracy but alter NC1, NC2, NC3, or NC0 can therefore
leave one detector family largely compatible while changing another.

\subsection{Diagnostic controls and empirical link}
\label{subsec:diagnostic-controls}

The preceding subsections identify the geometry channels through which
optimizer choice can affect feature-based OOD scores. The experiments use
diagnostic controls to localize these channels; they are not presented as
universal detector fixes.

First, we use post-hoc L2 normalization,
\begin{equation}
    \widehat f_o(x)
    =
    \frac{f_o(x)}
    {\|f_o(x)\|_2+\epsilon},
    \label{eq:posthoc-l2-normalization}
\end{equation}
and recompute distance- and neighborhood-based scores. Recovery after L2
normalization indicates that radial or feature-norm effects are an important
part of the raw detector gap. Lack of recovery indicates that angular,
covariance, local-density, or subspace structure remains detector-relevant.

Second, for Gaussian density readouts we compare covariance estimators. In
particular, the shrinkage covariance is
\begin{equation}
    \widehat{\Sigma}_{c,o,\alpha}
    =
    (1-\alpha)\widehat{\Sigma}_{c,o}
    +
    \alpha
    \frac{\operatorname{tr}(\widehat{\Sigma}_{c,o})}{d}
    I.
    \label{eq:covariance-shrinkage}
\end{equation}
Improvement under shrinkage indicates sensitivity to covariance conditioning or
finite-sample covariance estimation. Failure to recover under shrinkage
suggests that the detector gap cannot be explained by covariance estimation
alone.

Third, we measure geometry and detector changes in parallel. For geometry
metric \(g\), detector family \(m\), and OOD dataset or regime \(q\), the
experiments summarize changes relative to a reference optimizer as
\begin{equation}
    \Delta G_{o,g}
    =
    G_{o,g}
    -
    G_{o_{\mathrm{ref}},g},
    \qquad
    \Delta A_{o,m,q}
    =
    A_{o,m,q}
    -
    A_{o_{\mathrm{ref}},m,q}.
    \label{eq:delta-geometry-detector}
\end{equation}
Component-level correlations, profile associations, Mantel-style tests, and
geometry--detector association summaries are reported in
Section~\ref{sec:experiments} or the appendix as diagnostic evidence. They are
not used as causal proofs. Their role is to test whether the geometry channels
defined in this section move consistently with the detector-family behavior
observed empirically.
